% TV1 phụ trách -- Tổng quan & Kiến thức nền
% Thành viên: Trần Đình Luân (23120059)
 
% =====================================================================
\section{Tổng quan chủ đề}
% =====================================================================
 
\subsection{Giới thiệu bài toán}
 
Năm 1770, Wolfgang von Kempelen giới thiệu với châu Âu một cỗ máy
biết chơi cờ mang tên ``Mechanical Turk'' -- một thiết bị tinh vi với
hàng chục bánh răng và đòn bẩy. Nó đánh bại Napoleon, Benjamin
Franklin và vô số đối thủ khác. Hóa ra, toàn bộ bộ máy chỉ là vỏ
bọc, bên trong có một kỳ thủ thật ngồi điều khiển mọi nước đi qua hệ
thống gương phản chiếu.
 
Câu chuyện này phản ánh chính xác trạng thái của AI ngày nay. Phần
lớn các hệ thống AI hiện đại -- từ GPT-4, Gemini cho đến Stable
Diffusion -- đều được xây dựng trên nguyên tắc \textbf{maximum
likelihood imitation}: mô hình học cách tái tạo hành vi của con người
bằng cách tối đa hóa xác suất sinh ra dữ liệu giống dữ liệu huấn
luyện. Phương pháp này rất mạnh, nhưng chứa đựng một giới hạn nền
tảng không thể tránh khỏi:
 
\begin{equation}
    \text{Khả năng của mô hình} \;\leq\; \text{Khả năng của dữ liệu
    huấn luyện}
    \label{eq:imitation_bound}
\end{equation}
 
Nếu tất cả dữ liệu đều sai, mô hình sẽ sai theo. Nếu chưa có con
người nào giải được một bài toán, mô hình không thể học cách giải nó
từ dữ liệu đó. Đây là bài toán cốt lõi mà tutorial
\textit{Generative AI Meets Reinforcement Learning}
\cite{eysenbach2025generativerl} đặt ra: làm thế nào để xây dựng AI có
thể \textbf{vượt qua giới hạn của dữ liệu huấn luyện}?
 
Tutorial được trình bày tại ICML 2025 bởi Benjamin Eysenbach (Princeton
University) và Amy Zhang (UT Austin), xoay quanh hai câu hỏi liên
quan chặt chẽ:
\begin{itemize}
    \item Học Tăng cường và AI Tạo sinh -- hai lĩnh vực tưởng
    như tách biệt -- liệu có chia sẻ cùng cấu trúc toán học không?
    \item Nếu có, chúng ta có thể khai thác điều đó để xây dựng các hệ
    thống AI mạnh mẽ hơn như thế nào?
\end{itemize}
 
\subsection{Động lực nghiên cứu}
 
\subsubsection{Giới hạn của Học Bắt chước}
 
Như đã trình bày, phần lớn AI hiện đại được huấn luyện bằng cách tối
đa hóa khả năng tái tạo hành vi của con người. Cách tiếp cận này hiệu
quả khi có đủ dữ liệu chất lượng cao và nhiệm vụ nằm trong phân phối
dữ liệu. Tuy nhiên, nó gặp khó khăn cơ bản khi bài toán đòi hỏi sự
sáng tạo vượt ngoài dữ liệu, hoặc khi chưa có ai giải được bài toán
đó trước đây.
 
Một hệ quả đáng lo ngại hơn nữa là hiện tượng \textit{sụp đổ mô hình}:
khi AI được huấn luyện đệ quy trên dữ liệu do chính AI sinh ra, chất
lượng suy giảm theo cấp số nhân \cite{shumailov2024ai}.
 
\subsubsection{Giới hạn của Thiết kế Hàm Thưởng}
 
Nền tảng của Học Tăng cường là tín hiệu phần thưởng (reward).
Tuy nhiên, hầu hết bài toán thực tế rất khó thiết kế hàm thưởng
phù hợp -- đây là mâu thuẫn cốt lõi: Học Tăng cường hứa hẹn tìm ra giải pháp
tốt hơn con người, nhưng chính hàm thưởng lại đòi hỏi con người
phải hiểu rõ vấn đề mới có thể thiết kế được.
 
Với bài toán dọn phòng robot, người thiết kế phải tự hỏi: nhặt một
chiếc áo được bao nhiêu điểm? Xếp sách ngay ngắn được bao nhiêu? Nếu
thiết kế sai, robot có thể học cách ``gian lận'' thay vì thực sự dọn
dẹp -- một hiện tượng được gọi là \textit{khai thác phần thưởng}.
 
\subsubsection{Sự hội tụ của hai lĩnh vực}
 
Trong nhiều thập kỷ, Học Tăng cường và AI Tạo sinh phát triển song song nhưng
tương đối độc lập. Gần đây, giới nghiên cứu nhận ra rằng hai lĩnh vực
này đang giải quyết cùng một bài toán từ hai góc nhìn khác nhau, và
khai thác điểm chung này mở ra ba hướng giao thoa chính:
 
\begin{table}[H]
    \centering
    \caption{Ba hướng giao thoa giữa AI Tạo sinh và Học Tăng cường.}
    \label{tab:convergence}
    \begin{tabular}{ll}
        \toprule
        \textbf{Hướng} & \textbf{Ý nghĩa} \\
        \midrule
        AI Tạo sinh $\to$ Học Tăng cường   & Dùng mô hình tạo sinh làm thành phần trong Học Tăng cường \\
        Học Tăng cường $\to$ AI Tạo sinh   & Nhìn tương tác Học Tăng cường như một quá trình tạo sinh \\
        Học Tăng cường để huấn luyện AI Tạo sinh & Dùng Học Tăng cường để cải thiện cách huấn luyện mô hình tạo sinh \\
        \bottomrule
    \end{tabular}
\end{table}
 
\subsection{Tầm quan trọng và tính thời sự}
 
\subsubsection{Ứng dụng thực tiễn}
 
Sự kết hợp giữa AI Tạo sinh và Học Tăng cường có nhiều ứng dụng thực tiễn quan
trọng \cite{chi2023diffusion} \cite{ouyang2022instructgpt}:
 
\textbf{Robot học:} Diffusion policies đã đạt kết quả tốt nhất trong
nhiều tác vụ điều khiển robot. Kết hợp với Học Tăng cường, robot có thể vượt qua
giới hạn của dữ liệu minh họa, tự học cách thực hiện các hành vi
chưa từng được quan sát.
 
\textbf{Mô hình Nền tảng Hành vi:} Tương tự như GPT là mô hình nền tảng cho ngôn ngữ, các Mô hình Nền tảng Hành vi -- được học thông
qua Học Tăng cường không có phần thưởng -- có thể đóng vai trò mô hình nền tảng cho
hành động, một model dùng được cho nhiều tác vụ robot khác nhau mà
không cần huấn luyện lại.
 
\textbf{Học Tăng cường từ Phản hồi Con người và Căn chỉnh AI:} Học Tăng cường từ Phản hồi Con người
(RLHF) -- kỹ thuật cốt lõi làm cho ChatGPT ``nghe lời'' --
là ứng dụng trực tiếp của Học Tăng cường để tinh chỉnh mô hình tạo sinh. Hiểu
sâu mối liên hệ này giúp thiết kế các phương pháp căn chỉnh tốt hơn.
 
\textbf{Khoa học và khám phá:} AlphaFold dự đoán cấu trúc protein,
AlphaProof giải bài thi toán Olympic -- các hệ thống này đều dùng Học Tăng cường
để tìm giải pháp mà không có dữ liệu tương ứng từ trước.
 
\subsubsection{Tổng quát hóa -- bài toán còn bỏ ngỏ}
 
Tutorial đặt ra một câu hỏi quan trọng về tương lai: điều gì làm cho
mô hình tạo sinh hình ảnh và mô hình ngôn ngữ trở nên ấn tượng không phải
là chúng ghi nhớ dữ liệu đã thấy, mà là chúng \textit{tổng quát hóa} --
tạo ra output đúng trên những input chưa bao giờ thấy trong training.
 
Trong khi đó, các tác nhân Học Tăng cường ngày nay phần lớn vẫn chỉ ghi nhớ thay vì
tổng quát hóa. Tutorial đặt câu hỏi: tương tự như một mô hình tạo sinh hình ảnh
có thể tạo ra ảnh sống động sau khi chỉ thấy một phần nhỏ các
ảnh có thể, liệu ta có thể xây dựng các tác nhân Học Tăng cường có thể giải quyết các
bài toán mới sau khi chỉ practice trên một phần nhỏ không gian bài
toán đó? Đây chính là \textit{biên giới tổng quát hóa} mà kết hợp
AI Tạo sinh và Học Tăng cường hướng tới.
 
\subsection{Cấu trúc báo cáo}
 
Báo cáo được tổ chức thành các phần chính như sau:
\begin{itemize}
    \item \textbf{Phần 2} trình bày kiến thức nền về Học Tăng cường
    và Mô hình Tạo sinh, các khái niệm cơ sở cần thiết để
    hiểu nội dung tutorial.
    \item \textbf{Phần 3} trình bày nội dung chính của tutorial, bao
    gồm Phần I (Mô hình Tạo sinh trong Học Tăng cường), Phần II (Tương tác như
    một Mô hình Tạo sinh), Phần III (Học Hàm Khả năng) và Phần IV (Dữ liệu
    Tự Sinh).
    \item \textbf{Phần 4} mô tả thực nghiệm và kết quả demo nhóm thực
    hiện.
    \item \textbf{Phần 5} đánh giá ưu nhược điểm của các phương pháp
    và thảo luận hướng nghiên cứu tương lai.
\end{itemize}
 
% =====================================================================
\section{Kiến thức nền}
% =====================================================================
 
\subsection{Học Tăng cường}

% ---------------------------------------------------------------------
\subsubsection{Vòng lặp Tác nhân--Môi trường: Bức tranh tổng thể}
% ---------------------------------------------------------------------

Trước khi đi vào công thức toán học, hãy hình dung Học Tăng cường hoạt động như
thế nào theo trực giác. Trong Học Tăng cường, một \textbf{tác nhân} (ví dụ: robot,
chương trình chơi game, hệ thống điều khiển) liên tục tương tác với
một \textbf{môi trường} (\textit{environment}) theo vòng lặp lặp đi
lặp lại:

\begin{enumerate}
    \item Tác nhân quan sát \textbf{trạng thái hiện tại} $s_t$ của môi
    trường (ví dụ: hình ảnh màn hình game, vị trí các khớp robot).
    \item Tác nhân chọn một \textbf{hành động} $a_t$ dựa trên
    \textit{chính sách} (chiến lược) của mình (ví dụ: nhảy, đi thẳng,
    xoay trái).
    \item Môi trường chuyển sang \textbf{trạng thái mới} $s_{t+1}$
    theo xác suất $p(s_{t+1} \mid s_t, a_t)$, đồng thời phát ra
    \textbf{tín hiệu phần thưởng} $r_t = r(s_t, a_t)$ (ví dụ: $+1$
    mỗi khi xóa được một dòng trong Tetris, $-1$ khi robot ngã).
    \item Lặp lại từ bước 1.
\end{enumerate}

\begin{figure}[H]
    \centering
    \begin{tikzpicture}[
        box/.style={draw, rounded corners=3pt, align=center,
                    minimum width=2.8cm, minimum height=1.0cm,
                    fill=white},
        arr/.style={-{Latex[length=2.5mm]}, thick}
    ]
        \node[box] (agent) at (0, 0)   {Tác nhân\\$\pi(a \mid s)$};
        \node[box] (env)   at (6, 0)   {Môi trường\\$p(s' \mid s, a)$};

        % action: agent -> env (top arc)
        \draw[arr] (agent.north) to[out=60, in=120]
            node[midway, above]{hành động $a_t$} (env.north);

        % state + reward: env -> agent (bottom arc)
        \draw[arr] (env.south) to[out=-120, in=-60]
            node[midway, below]{trạng thái $s_{t+1}$, phần thưởng $r_t$}
            (agent.south);
    \end{tikzpicture}
    \caption{Vòng lặp tác nhân--môi trường trong Học Tăng cường. Tác nhân gửi hành động
    cho môi trường; môi trường phản hồi trạng thái mới và phần thưởng.
    Mục tiêu tác nhân là học chính sách $\pi$ tối đa hóa tổng phần thưởng
    tích lũy.}
    \label{fig:rl-loop}
\end{figure}

Điểm khác biệt căn bản so với Học Có Giám sát là: \textbf{không
có ``đáp án đúng'' nào được cung cấp sẵn}. Tác nhân phải tự khám phá
thông qua thử-và-sai, và tín hiệu học duy nhất là phần thưởng thưa
thớt từ môi trường. Điều này đồng thời là điểm mạnh (tác nhân có thể
tìm ra chiến lược vượt trội hơn bất kỳ dữ liệu huấn luyện nào) và là
thách thức chính của Học Tăng cường.

% ---------------------------------------------------------------------
\subsubsection{Định nghĩa bài toán Học Tăng cường: Quá trình Quyết định Markov}
% ---------------------------------------------------------------------
 
Bài toán Học Tăng cường được hình thức hóa bằng khung
\textbf{Quá trình Quyết định Markov (MDP)}, bao gồm năm thành phần:
 
\begin{equation}
    \text{MDP} = (\mathcal{S},\; \mathcal{A},\; p,\; r,\; \gamma)
    \label{eq:mdp}
\end{equation}
 
\begin{table}[H]
    \centering
    \caption{Các thành phần của Quá trình Quyết định Markov, minh họa qua game Tetris.}
    \label{tab:mdp}
    \begin{tabular}{lll}
        \toprule
        \textbf{Ký hiệu} & \textbf{Tên} & \textbf{Ví dụ (Tetris)} \\
        \midrule
        $\mathcal{S}$ & Không gian trạng thái & Mọi cấu hình bảng có thể \\
        $\mathcal{A}$ & Không gian hành động & \{trái, phải, xoay, thả\} \\
        $p(s'\mid s,a)$ & Hàm chuyển trạng thái & Xác suất khối rơi vào ô nào \\
        $r(s,a)$ & Hàm phần thưởng & $+1$ mỗi dòng xóa được \\
        $\gamma \in [0,1)$ & Hệ số chiết khấu & Thường từ $0.9$ đến $0.999$ \\
        \bottomrule
    \end{tabular}
\end{table}
 
Giả định \textbf{Markov} là nền tảng quan trọng của khung này: trạng
thái hiện tại $s_t$ chứa đủ thông tin để quyết định hành động tiếp
theo, không cần nhớ toàn bộ lịch sử. Nói cách khác, ``quá khứ'' chỉ
ảnh hưởng đến tương lai thông qua trạng thái hiện tại. Nhờ tính Markov,
ta có thể dùng các phương trình đệ quy (phương trình Bellman) để tính
toán hiệu quả thay vì phải xét toàn bộ lịch sử.
 
\subsubsection{Chính sách và Mục tiêu Học Tăng cường}
 
\textbf{Chính sách} $\pi$ là chiến lược ra quyết định của tác nhân -- một
hàm ánh xạ từ trạng thái sang hành động (hoặc phân phối trên hành
động):
\[
    \pi(a \mid s) = \text{xác suất chọn hành động } a \text{ khi ở trạng thái }
    s
\]

Ví dụ: chính sách của người chơi Tetris giỏi sẽ gán xác suất cao cho
hành động ``xoay'' khi nhìn thấy khối L ở vị trí nhất định.

Một \textbf{quỹ đạo} $\tau = (s_0, a_0, s_1, a_1, \ldots)$
là chuỗi trạng thái và hành động theo thời gian -- về bản chất là một
``ván chơi'' hoàn chỉnh. \textbf{Lợi tức chiết khấu} $R(\tau)$ là
tổng phần thưởng có chiết khấu theo thời gian:
 
\begin{equation}
    R(\tau) = \sum_{t=0}^{\infty} \gamma^t \, r(s_t, a_t)
    \label{eq:return}
\end{equation}
 
Hệ số chiết khấu $\gamma < 1$ giải quyết hai vấn đề: nó đảm bảo tổng
phần thưởng hội tụ (không vô hạn), đồng thời phản ánh nguyên tắc kinh
tế học rằng phần thưởng gần hơn có giá trị hơn (giống như tiền ngày
hôm nay có giá trị hơn tiền ngày mai).
 
\textbf{Mục tiêu Học Tăng cường} là tìm chính sách $\pi$ tối đa hóa lợi tức kỳ vọng --
tức phần thưởng tích lũy kỳ vọng trên tất cả các quỹ đạo có thể
xảy ra:
 
\begin{equation}
    J(\pi) = \mathbb{E}_{\tau \sim \pi} \left[ R(\tau) \right]
    = \mathbb{E}_{\tau \sim \pi} \left[ \sum_{t=0}^{\infty} \gamma^t
    r(s_t, a_t) \right]
    \label{eq:rl_objective}
\end{equation}

Ký hiệu $\tau \sim \pi$ nghĩa là quỹ đạo được sinh ra bởi chính
chính sách $\pi$ tương tác với môi trường: tại mỗi bước $t$, tác nhân chọn
$a_t \sim \pi(\cdot \mid s_t)$, môi trường phản hồi
$s_{t+1} \sim p(\cdot \mid s_t, a_t)$.
 
\subsubsection{So sánh Học Tăng cường và Học Có Giám sát}
 
Sự khác biệt căn bản giữa Học Tăng cường (HTc) và Học Có Giám sát (HGs) nằm ở nguồn
gốc dữ liệu:
 
\begin{equation}
    \text{HGs:} \quad \max_\theta \; \mathbb{E}_{(s,a) \sim \mathcal{D}}
    \bigl[\log \pi_\theta(a \mid s)\bigr]
    \label{eq:sl_obj}
\end{equation}
\begin{equation}
    \text{HTc:} \quad \max_\theta \; \mathbb{E}_{\tau \sim \pi_\theta}
    \bigl[R(\tau)\bigr]
    \label{eq:rl_obj_compare}
\end{equation}
 
Trong HGs, dữ liệu đến từ tập dữ liệu cố định $\mathcal{D}$ và không phụ
thuộc vào tham số đang học. Trong HTc, quỹ đạo được sinh ra bởi
chính chính sách $\pi_\theta$ đang được tối ưu -- khi $\theta$ thay đổi,
cả phân phối hành động lẫn phân phối trạng thái được ghé thăm đều thay
đổi theo. Đây chính là điểm mạnh (tác nhân có thể khám phá) nhưng cũng
là thách thức chính (dữ liệu \textit{không dừng}) của Học Tăng cường.
 
Ví dụ trực quan: nếu dùng HGs để học chơi Tetris trên tập dữ liệu từ người
chơi mới bắt đầu, model sẽ không thể giỏi hơn người đó. Nhưng với HTc,
tác nhân tự tạo dữ liệu thông qua tương tác và có thể tìm ra chiến lược
tốt hơn bất kỳ dữ liệu nào trong tập dữ liệu \cite{eysenbach2025generativerl}.
 
\subsubsection{Hàm Giá trị và Phương trình Bellman}
 
Để tối ưu $J(\pi)$, cần biết ``từ trạng thái này, hành động nào dẫn
đến lợi tức cao nhất?'' -- đây là vai trò của \textbf{hàm giá trị}.
 
\textbf{Hàm Q} $Q^\pi(s, a)$ cho biết lợi tức kỳ vọng khi bắt đầu
từ cặp $(s, a)$ và sau đó theo chính sách $\pi$:
 
\begin{equation}
    Q^\pi(s, a) = \mathbb{E}_\pi \left[ \sum_{t=0}^{\infty} \gamma^t
    r(s_t, a_t) \;\Big|\; s_0 = s,\, a_0 = a \right]
    \label{eq:q_function}
\end{equation}

Diễn giải đơn giản: $Q^\pi(s, a)$ trả lời câu hỏi ``Nếu tôi đang ở
trạng thái $s$, thực hiện hành động $a$, rồi từ đó trở đi luôn theo
chính sách $\pi$, tôi kỳ vọng nhận được bao nhiêu phần thưởng tổng
cộng?''
 
Hàm Q thỏa \textbf{phương trình Bellman} -- quan hệ đệ quy cơ bản
của Học Tăng cường. Phương trình này phát biểu rằng giá trị hiện tại bằng phần
thưởng ngay lập tức cộng với giá trị tương lai được chiết khấu:
 
\begin{equation}
    Q^\pi(s, a) = \underbrace{r(s, a)}_{\text{thưởng ngay}} +
    \gamma \, \mathbb{E}_{s' \sim p(\cdot|s,a),\;
    a' \sim \pi(\cdot|s')} \bigl[ \underbrace{Q^\pi(s', a')}_{\text{giá trị tương lai}} \bigr]
    \label{eq:bellman}
\end{equation}
 
Phương trình Bellman quan trọng vì nó biến bài toán tối ưu hóa vô hạn
chiều thành bài toán có thể giải lặp đi lặp lại (quy hoạch động).
Hàm Q xuất hiện xuyên suốt tutorial -- trong chính sách khuếch tán
(dùng $Q$ để chọn hành động), trong học tương phản thời gian
($Q \approx$ tỉ lệ log-likelihood), và trong Học Tăng cường không cần huấn luyện lại (tính $Q$
từ biểu diễn học được).
 
\subsubsection{Gradient Chính sách}
 
Để tối đa hóa $J(\pi_\theta)$, ta lấy gradient theo tham số $\theta$
và thực hiện leo dốc gradient. Định lý Gradient Chính sách cho kết quả:
 
\begin{equation}
    \nabla_\theta J(\pi_\theta) = \mathbb{E}_{\tau \sim \pi_\theta}
    \left[ R(\tau) \cdot \sum_t \nabla_\theta \log \pi_\theta(a_t \mid
    s_t) \right]
    \label{eq:policy_gradient}
\end{equation}

Trực giác của gradient chính sách rất tự nhiên: số hạng
$\nabla_\theta \log \pi_\theta(a \mid s)$ đẩy tham số theo hướng tăng
xác suất của hành động $a$ tại trạng thái $s$; nhân với $R(\tau)$
nghĩa là ta tăng xác suất của những hành động đã dẫn đến lợi tức cao,
và giảm xác suất của những hành động dẫn đến lợi tức thấp.
 
Số hạng $\nabla_\theta \log \pi_\theta(a \mid s)$ được gọi là
\textbf{hàm điểm số} của chính sách. Đây là cầu nối kỹ thuật quan trọng
nhất giữa Học Tăng cường và mô hình tạo sinh: trong các mô hình khuếch tán dựa trên điểm số,
đại lượng tương tự $\nabla_x \log p(x)$ được dùng để hướng
dẫn quá trình sinh mẫu. Hai bài toán -- tối ưu chính sách và sinh mẫu từ
mô hình tạo sinh -- đều có thể được giải qua cùng một công cụ toán
học: \textbf{theo dõi gradient của log-xác suất}.

% ---------------------------------------------------------------------
\subsection{Mô hình Tạo sinh}
% ---------------------------------------------------------------------
 
\subsubsection{Tổng quan và Phân loại}
 
\textbf{Mô hình tạo sinh} là mô hình có khả năng (1) tạo ra mẫu mới
từ một phân phối học được, và (2) ước lượng hàm khả năng (xác suất) của
các mẫu đã cho. Khác với các mô hình phân biệt (phân loại đầu vào),
mô hình tạo sinh học phân phối xác suất $p(x)$ của dữ liệu và có thể
sinh ra mẫu mới $x \sim p(x)$.

Ý tưởng chung của các mô hình tạo sinh mạnh nhất hiện nay là thực
hiện \textbf{tính toán lặp đi lặp lại} để tạo ra mẫu:
\begin{itemize}
    \item \textbf{VAE} và luồng chuẩn hóa biến đổi vector nhiễu qua
    nhiều lớp mạng neural.
    \item \textbf{Mô hình khuếch tán} lặp đi lặp lại quá trình khử nhiễu
    qua hàng trăm bước.
    \item \textbf{Transformer} sinh token từng bước một theo thứ tự.
    \item \textbf{Mô hình khớp luồng} tích phân một trường vector để
    biến đổi phân phối.
\end{itemize}

Đặc điểm ``nhiều bước'' này không phải ngẫu nhiên -- nó phản ánh một
nguyên lý sâu xa: các phân phối phức tạp rất khó nắm bắt trong một
bước duy nhất, nhưng lại có thể được xây dựng dần dần qua nhiều bước
đơn giản hơn. Tutorial chỉ ra rằng \textbf{đây chính là điểm tương
đồng với Học Tăng cường}: một tác nhân Học Tăng cường cũng tạo ra kết quả thông qua nhiều bước
tương tác với môi trường.

\subsubsection{Bộ Mã hóa Biến phân (VAE)}

\textbf{VAE} là một trong những mô hình tạo sinh đầu tiên học được
không gian tiềm ẩn (latent space) ý nghĩa. VAE gồm hai thành phần:
bộ mã hóa $q_\phi(z \mid x)$ nén dữ liệu $x$ vào vector tiềm ẩn $z$,
và bộ giải mã $p_\theta(x \mid z)$ tái tạo dữ liệu từ $z$. Quá trình
sinh mẫu mới đơn giản là lấy mẫu $z \sim \mathcal{N}(0, I)$ rồi
cho qua bộ giải mã.

VAE được huấn luyện bằng cách tối đa hóa ELBO (Cận dưới Bằng chứng):
\begin{equation}
    \mathcal{L}_{\text{VAE}} = \mathbb{E}_{q_\phi(z|x)}
    [\log p_\theta(x \mid z)]
    - D_{\text{KL}}\bigl(q_\phi(z \mid x) \;\Vert\; \mathcal{N}(0,I)\bigr)
    \label{eq:vae_elbo}
\end{equation}

Số hạng đầu khuyến khích bộ giải mã tái tạo tốt; số hạng KL
chính quy hóa không gian tiềm ẩn về phân phối Gaussian chuẩn.

\textbf{Kết nối với Học Tăng cường:} Trong học kỹ năng (Phần IV của tutorial),
biến tiềm ẩn $z$ của VAE tương đồng chính xác với ``nhãn kỹ năng'' $z$
-- một biến điều kiện quyết định loại hành vi tác nhân thực hiện. Cả hai
đều nén thông tin về một tập hành vi phức tạp vào một không gian nhỏ
hơn có thể tìm kiếm được.

\subsubsection{Mạng Đối kháng Tạo sinh (GAN)}

\textbf{GAN} gồm hai mạng thi đấu với nhau theo dạng trò chơi
minimax: \textit{bộ tạo sinh} $G$ cố gắng sinh mẫu giống thật,
\textit{bộ phân biệt} $D$ cố gắng phân biệt mẫu thật với mẫu giả:
\begin{equation}
    \min_G \max_D \; \mathbb{E}_{x \sim p_{\text{data}}}[\log D(x)]
    + \mathbb{E}_{z \sim \mathcal{N}(0,I)}[\log(1 - D(G(z)))]
    \label{eq:gan}
\end{equation}

Khi đạt cân bằng Nash, bộ tạo sinh học được phân phối dữ liệu thật.

\textbf{Kết nối với Học Tăng cường:} Bộ phân biệt của GAN có thể được tái sử dụng
trực tiếp làm \textbf{tín hiệu phần thưởng} cho Học Tăng cường. Nếu bộ phân biệt học
được ``mẫu này trông như thật không'', ta có thể dùng $\log D(x)$ làm
phần thưởng để dạy tác nhân thực hiện hành vi trông tự nhiên -- đây là một
trong những cách phá vỡ ``rào cản trừu tượng hóa'' giữa mô hình tạo sinh
và thuật toán Học Tăng cường mà tutorial nhấn mạnh.

\subsubsection{Mô hình Khuếch tán}
 
Mô hình khuếch tán là họ mô hình tạo sinh hiện đại nhất, đứng sau các
hệ thống tạo ảnh mạnh nhất hiện nay như Stable Diffusion và DALL-E.
Ý tưởng cốt lõi: \textbf{học cách đảo ngược một quá trình phá hủy}.

\paragraph{Quá trình tiến (thêm nhiễu).}

Bắt đầu từ dữ liệu sạch $x_0$, ta dần dần thêm nhiễu Gaussian qua
$T$ bước để thu được chuỗi $x_1, x_2, \ldots, x_T$, trong đó $x_T$
gần với nhiễu trắng hoàn toàn:
\[
    q(x_t \mid x_{t-1}) = \mathcal{N}\bigl(x_t;\, \sqrt{1-\beta_t}\,
    x_{t-1},\; \beta_t \mathbf{I}\bigr)
\]
Nhờ tính chất của Gaussian, ta có thể tính thẳng $x_t$ từ $x_0$:
\[
    q(x_t \mid x_0) = \mathcal{N}(x_t;\, \sqrt{\bar{\alpha}_t}\,x_0,\;
    (1-\bar{\alpha}_t)\mathbf{I}), \quad
    \bar{\alpha}_t = \prod_{s=1}^t (1-\beta_s)
\]

\paragraph{Quá trình ngược (khử nhiễu).}

Mô hình học cách đảo ngược quá trình trên -- từ $x_T$ nhiễu, dần
dần tái tạo dữ liệu sạch $x_0$ qua nhiều bước. Đây là bước khó vì
cần ước lượng $p(x_{t-1} \mid x_t)$. Kỹ thuật mấu chốt là học
\textbf{hàm điểm số}:
\[
    s_\theta(x_t, t) \approx \nabla_{x_t} \log p_t(x_t)
\]
Hàm điểm số $\nabla_{x_t} \log p_t(x_t)$ là gradient của
log-xác suất theo $x_t$ -- nó chỉ hướng ``đi về phía vùng mật độ
xác suất cao hơn'' trong không gian dữ liệu tại bước thời gian $t$.

\textbf{Kết nối với Học Tăng cường:} Hàm điểm số của mô hình khuếch tán
$\nabla_x \log p_t(x)$ và hàm điểm số của gradient chính sách
$\nabla_\theta \log \pi_\theta(a \mid s)$ là cùng một cấu trúc toán
học -- gradient của log-xác suất theo biến cần tối ưu. Đây là lý
do mô hình khuếch tán và gradient chính sách có thể được tích hợp tự nhiên
vào nhau, như được minh họa qua DDPO và Diffusion-DICE ở Phần III.

Trong Học Tăng cường, mô hình khuếch tán được dùng làm \textbf{chính sách khuếch tán}:
thay vì chính sách Gaussian đơn giản chỉ biểu diễn một đỉnh, tác nhân dùng
một mô hình khuếch tán để biểu diễn phân phối hành động phức tạp, đa
đỉnh -- quan trọng khi tập dữ liệu chứa nhiều chiến lược khác nhau.

\subsubsection{Mô hình Khớp Luồng}
 
Khớp luồng là một hướng tiếp cận thay thế cho khuếch tán, học cách
biến đổi trực tiếp một phân phối đơn giản (thường là Gaussian) thành
phân phối dữ liệu thông qua một \textbf{trường vector} $v_\theta(x, t)$
có thể học được. Quá trình biến đổi được mô tả bởi phương trình vi
phân thường (ODE):
\begin{equation}
    \frac{dx}{dt} = v_\theta(x, t), \quad x(0) \sim \mathcal{N}(0, I),
    \quad x(1) \sim p_{\text{data}}
    \label{eq:flow_ode}
\end{equation}

Ý tưởng trực quan: ta ``vẽ'' một đường đi trơn trong không gian xác
suất, từ điểm xuất phát (nhiễu Gaussian) đến điểm đích (dữ liệu thật).
Trường vector $v_\theta$ chỉ hướng và tốc độ di chuyển tại mỗi điểm.

\paragraph{Ưu điểm so với khuếch tán.}

Mô hình khuếch tán dùng lịch trình nhiễu cố định với hàng trăm bước nhỏ.
Khớp luồng học đường đi \textbf{thẳng hơn} trong không gian xác
suất, cho phép sinh mẫu với ít bước hơn (đôi khi chỉ 1 bước ODE).
Về lý thuyết, khớp luồng tối ưu hóa trực tiếp hàm khả năng của
quỹ đạo thay vì phải thông qua các bước gián tiếp.

\textbf{Kết nối với Học Tăng cường:} Trong tutorial, Luồng Hiệu chỉnh Thời gian
\cite{farebrother2025tdflows} áp dụng khớp luồng để ước lượng
\textit{hàm khả năng của trạng thái tương lai} -- tức thước đo kế tiếp
-- kết hợp với cấu trúc Bellman để tận dụng dữ liệu ngoài chính sách. Đây
là ví dụ điển hình về việc khai thác đặc thù của một họ mô hình tạo sinh
để giải quyết bài toán Học Tăng cường.
 
\subsubsection{Transformer và Mô hình Tự Hồi quy}
 
Transformer là kiến trúc mạng neural dựa trên cơ chế
\textit{tự chú ý}, hiện đang chi phối phần lớn AI hiện đại (GPT,
BERT, ViT). Trước khi giải thích ứng dụng trong Học Tăng cường, cần hiểu hai khái
niệm cốt lõi:

\paragraph{Tự chú ý.}

Với một chuỗi token $(x_1, \ldots, x_n)$, tự chú ý tính
\textbf{mức độ liên quan} giữa mỗi cặp token và tổng hợp thông tin
theo mức độ đó:
\[
    \text{Attention}(Q, K, V) = \text{softmax}\!\left(
    \frac{QK^\top}{\sqrt{d_k}}\right) V
\]
trong đó $Q$ (truy vấn), $K$ (khóa), $V$ (giá trị) là các biến đổi tuyến
tính của đầu vào. Cơ chế này cho phép mỗi token ``chú ý'' đến những
token liên quan nhất, bất kể chúng cách nhau bao xa trong chuỗi --
điều mà RNN và LSTM gặp khó khăn.

\paragraph{Mô hình tự hồi quy.}

Mô hình tự hồi quy sinh chuỗi token từng bước một, trong đó mỗi
token được sinh có điều kiện trên toàn bộ các token đã sinh trước đó:
\[
    p(x_1, x_2, \ldots, x_T) = \prod_{t=1}^{T} p(x_t \mid x_1, \ldots,
    x_{t-1})
\]
Đây là nguyên tắc hoạt động của GPT: sinh token tiếp theo dựa trên
tất cả token trước đó.

\paragraph{Transformer Quyết định: mô hình tự hồi quy làm chính sách.}

\textbf{Transformer Quyết định} \cite{chen2021decision} áp dụng trực
tiếp kiến trúc này cho Học Tăng cường: thay vì học hàm giá trị hay thực hiện
cập nhật Bellman, mô hình coi một quỹ đạo Học Tăng cường -- chuỗi
$(R_0, s_0, a_0, R_1, s_1, a_1, \ldots)$ -- như một chuỗi token và
học sinh hành động có điều kiện trên \textit{lợi tức-còn-lại} (tổng
phần thưởng kỳ vọng từ bước hiện tại trở đi) và lịch sử quỹ đạo.

Ý tưởng then chốt: nếu ta điều kiện hóa mô hình trên ``lợi tức mục tiêu
cao'' (ví dụ: 100 điểm), model sẽ sinh ra những hành động từng được
quan sát trong các quỹ đạo có lợi tức cao. Đây là ví dụ điển hình
của việc dùng mô hình tạo sinh làm chính sách trong Học Tăng cường.

% ---------------------------------------------------------------------
\subsection{Thước đo Kế tiếp}
% ---------------------------------------------------------------------
 
\textbf{Thước đo kế tiếp} (hay thước đo chiếm dụng trạng thái chiết khấu)
$p^\pi(s_f)$ là một trong những khái niệm xuyên suốt toàn bộ tutorial.
Nó được định nghĩa là xác suất mà tác nhân đến được trạng thái $s_f$ tại
một thời điểm nào đó trong tương lai khi theo chính sách $\pi$:
 
\begin{equation}
    p^\pi(s_f) \triangleq (1-\gamma) \, \mathbb{E}_\pi \left[
    \sum_{t=0}^{\infty} \gamma^t \, p(s_{t+1} = s_f \mid s_t, a_t)
    \right]
    \label{eq:successor_measure_intro}
\end{equation}
 
Hệ số $(1-\gamma)$ chuẩn hóa để $p^\pi$ là một phân phối xác suất hợp
lệ (tổng bằng 1). Phiên bản có điều kiện $p^\pi(s_f \mid s, a)$ cho
biết xác suất đến $s_f$ khi bắt đầu từ trạng thái $s$ và thực hiện
hành động $a$.
 
\subsubsection{Ý nghĩa trực quan và ví dụ minh họa}
 
Thước đo kế tiếp trả lời câu hỏi: ``Nếu tác nhân bắt đầu từ trạng thái
$s$ và theo chính sách $\pi$, nó có bao nhiêu khả năng ghé thăm trạng thái
$s_f$ trong tương lai?'' Khái niệm này \textbf{tổng quát hóa
hàm Q}: trong khi hàm Q tính \textit{một số} (tổng phần
thưởng kỳ vọng), thước đo kế tiếp tính \textit{một phân phối đầy đủ}
trên các trạng thái tương lai.

\paragraph{Ví dụ minh họa bằng mê cung $3 \times 3$.}

Hãy tưởng tượng robot đứng tại ô $(1,1)$ trong mê cung $3\times 3$,
với chính sách ``luôn đi thẳng về phía đông''. Thước đo kế tiếp
$p^\pi(s_f \mid (1,1), \text{đi đông})$ với $\gamma = 0.9$ sẽ trông
như sau (giá trị xấp xỉ):

\begin{center}
\begin{tabular}{|c|c|c|}
    \hline
    ô $(1,1)$ & \textbf{ô $(1,2)$: 0.36} & \textbf{ô $(1,3)$: 0.32} \\
    & (sau 1 bước: $0.9^0 \cdot 0.4$) & (sau 2 bước) \\
    \hline
    \multicolumn{3}{|c|}{... xác suất thấp dần ở các ô phía sau ...} \\
    \hline
\end{tabular}
\end{center}

Khác với hàm Q chỉ cho biết ``tổng reward kỳ vọng là bao nhiêu'',
thước đo kế tiếp cho biết \textit{tác nhân sẽ ở đâu} -- thông tin phong
phú hơn. Lợi ích then chốt: nếu phần thưởng thay đổi (ví dụ: bây giờ ô
$(1,3)$ có thưởng thay vì ô $(1,2)$), ta chỉ cần tính lại tích phân
trên thước đo kế tiếp đã học, \textbf{không cần train lại từ đầu}.

\subsubsection{Liên hệ với Hàm Q}
 
Kết nối giữa thước đo kế tiếp và hàm Q rất trực tiếp. Nếu có
thước đo kế tiếp $p^\pi(s_f \mid s, a)$ và tập trạng thái có nhãn
thưởng $r(s_f)$, ta có thể tính hàm Q bằng:
 
\begin{equation}
    Q^\pi(s, a) = \mathbb{E}_{s_f \sim p^\pi(\cdot|s,a)} \left[
    r(s_f) \right] = \int r(s_f) \, p^\pi(s_f \mid s, a) \, ds_f
    \label{eq:q_from_succ}
\end{equation}
 
Điều này có ý nghĩa lớn: nếu học được thước đo kế tiếp một lần, ta có
thể tính hàm Q cho \textbf{bất kỳ hàm thưởng nào} mà không
cần huấn luyện lại -- đây là nền tảng cho Học Tăng cường không cần huấn luyện lại được trình bày
trong Phần III.
 
\subsubsection{Tính toán qua Bellman}
 
Thước đo kế tiếp thỏa mãn một phương trình Bellman tương tự
hàm Q:
 
\begin{equation}
    p^\pi(s_f \mid s, a) = (1-\gamma) \, p(s' = s_f \mid s, a)
    + \gamma \, \mathbb{E}_{s' \sim p(\cdot|s,a), \; a' \sim \pi(\cdot|s')}
    \bigl[ p^\pi(s_f \mid s', a') \bigr]
    \label{eq:bellman_succ}
\end{equation}

Diễn giải: xác suất đến $s_f$ từ $(s,a)$ bằng xác suất đến $s_f$
\textit{ngay bước tiếp theo} (số hạng đầu) cộng với xác suất đến $s_f$
\textit{sau nhiều bước} thông qua trạng thái trung gian $s'$ (số hạng
sau).

Tính chất này cho phép ước lượng thước đo kế tiếp từ dữ liệu
ngoài chính sách bằng học hiệu chỉnh thời gian, tương tự như cách ta ước
lượng hàm Q. Đây là lý do thước đo kế tiếp có thể được học hiệu
quả từ dữ liệu tương tác.
 
\subsubsection{Vị trí trong Tutorial}
 
Thước đo kế tiếp xuất hiện ở nhiều phần khác nhau của tutorial với
các vai trò đa dạng. Trong Phần II, nó mô tả phân phối trạng thái đích
của bài toán đạt đích. Trong Phần III, nó là hàm khả năng của
``mô hình tạo sinh tương tác'' cần được ước lượng. Trong Học Tăng cường không cần huấn luyện lại,
các biểu diễn tiến-lùi học cách phân tích thước đo kế tiếp để cho phép tính chính sách tối ưu cho bất kỳ hàm thưởng nào
chỉ trong vài bước tính toán.
 
% ---------------------------------------------------------------------
\subsection{Bảng đối chiếu Học Tăng cường và Mô hình Tạo sinh}
% ---------------------------------------------------------------------
 
Một trong những đóng góp quan trọng nhất của tutorial là chỉ ra sự
tương đương toán học giữa Học Tăng cường và Mô hình Tạo sinh. Bảng~\ref{tab:rl_genai}
tổng hợp các khái niệm tương ứng giữa hai lĩnh vực. Mỗi hàng trong
bảng thể hiện một cặp khái niệm đóng \textit{cùng một vai trò toán
học} trong hai lĩnh vực khác nhau -- hiểu bảng này giúp giải thích
tại sao kỹ thuật từ lĩnh vực này có thể áp dụng cho lĩnh vực kia.

Một số cặp chưa quen thuộc sẽ được giải thích trong Phần III:
\textit{thước đo chiếm dụng} $d^\pi(s,a)$ là phân phối tần suất
trạng thái-hành động mà chính sách ghé thăm (tổng quát hóa thước đo kế tiếp sang
không gian trạng thái-hành động); \textit{tỉ lệ hàm khả năng} là tỉ số
$p^\pi(s_f \mid s,a) / p(s_f)$ đo mức độ chính sách tăng xác suất đến
$s_f$ so với đường cơ sở.
 
\begin{table}[H]
    \centering
    \caption{Các khái niệm tương đương giữa Học Tăng cường và Mô hình Tạo sinh.}
    \label{tab:rl_genai}
    \begin{tabular}{lll}
        \toprule
        \textbf{Học Tăng cường} & \textbf{Mô hình Tạo sinh} &
        \textbf{Ý nghĩa chung} \\
        \midrule
        Chính sách $\pi(a \mid s)$ & Mô hình tạo sinh $p(x)$ & Phân phối
        cần học \\
        Hàm điểm số $\nabla_\theta \log \pi(a \mid s)$ & Hàm điểm số
        $\nabla_x \log p(x)$ & Gradient của log-xác suất \\
        Lợi tức $R(\tau)$ & Log-hàm khả năng $\log p(x)$ & Tín hiệu học \\
        Quỹ đạo $\tau$ & Biến tiềm ẩn $z$ & Biến tiềm ẩn \\
        Gradient chính sách & Ước lượng hàm điểm số & Cách tối ưu phân
        phối \\
        Hàm Q $Q(s,a)$ & Tỉ lệ hàm khả năng
        $\frac{p^\pi(s_f|s,a)}{p(s_f)}$ & Ước lượng xác suất tương
        lai \\
        Thước đo chiếm dụng $d^\pi(s,a)$ & Phân phối dữ liệu $p(x)$ &
        Phân phối muốn khớp \\
        Phương trình Bellman & Tính chất Markov & Cấu trúc đệ quy \\
        Nhãn kỹ năng $z$ & Mã tiềm ẩn $z$ trong VAE & Điều khiển hành
        vi \\
        \bottomrule
    \end{tabular}
\end{table}
 
Hiểu bảng này giúp đọc phần nội dung chính dễ hơn đáng kể: mỗi khi
tutorial dùng kỹ thuật từ một lĩnh vực để giải quyết bài toán của
lĩnh vực kia, đều có thể truy về một cặp tương đương trong
bảng~\ref{tab:rl_genai}.